\subsection{Integration Kernels for 3D Observables}\label{subsec:obs_sfb_kernels}

This section develops the integration kernels that connect underlying matter power spectra to observable weak lensing and clustering.

\subsubsection{Generic kernel structure}

Observable power spectra take the form:

\begin{equation}
C_\ell = \int_0^\infty \frac{dk}{2\pi^2} k^2 P(k) I^2_\ell(k)
\end{equation}

where $I_\ell(k)$ is an integration kernel that encodes the survey geometry and physics.

\subsubsection{Convergence kernel}

For weak lensing convergence:

\begin{equation}
I^{\kappa}_\ell(k) = \int_0^\infty d\chi \, W^\kappa(\chi) j_\ell(k\chi)
\end{equation}

where $W^\kappa(\chi)$ is the convergence lensing efficiency (depends on source redshift distribution and angular diameter distances).

\subsubsection{Density/clustering kernel}

For galaxy clustering:

\begin{equation}
I^{\delta}_\ell(k) = \int_0^\infty d\chi \, n(\chi) b(\chi) D_+(\chi) j_\ell(k\chi)
\end{equation}

where $n(\chi)$ is galaxy number density, $b(\chi)$ is bias, and $D_+(\chi)$ is growth factor.

\subsubsection{Numerical evaluation}

For accurate predictions:
1. Compute lensing/clustering weights: $W(\chi)$ or $n(\chi) b(\chi) D_+(\chi)$
2. Evaluate Bessel functions $j_\ell(k\chi)$ using recurrence or SpecialFunctions.jl
3. Integrate over $\chi$ using adaptive quadrature (QuadGK.jl)
4. Convolve with power spectrum

See \cref{ch:numerical} for integration algorithms.
